\input{../tex-inputs/latex-korrekturansicht-vorspann}

\section[Arthur Schnitzler an Gustav Schwarzkopf, 31. 8. 1924]{L04355 Arthur Schnitzler an Gustav Schwarzkopf, 31. 8. 1924}
\nopagebreak\mylabel{L04355v}
\rehead{ }\normalsize\beginnumbering\briefempfaengerindex{Schwarzkopf, Gustav@\textsc{Schwarzkopf, Gustav}!zzzSchnitzler, Arthur@\emph{von Arthur Schnitzler}!1924-08-311@{31. 8. 1924}|(be}
\toendnotes[C]{\smallbreak\pagebreak[2]}
\correspDesc{Versand  durch Arthur Schnitzler am 31. 8. 1924 in Locarno
\newline{}Übermittlung  am 31. 8. 1924 in Bellinzona
\newline{}Erhalt  durch Gustav Schwarzkopf im Zeitraum [1. 9. 1924 – 5. 9. 1924?] in Wien}\toendnotes[C]{\smallbreak}
\Standort{DLA, A:Schnitzler, HS.1985.1.1897.}
\physDesc{Bildpostkarte, 240 Zeichen
\newline{}Handschrift: Bleistift, lateinische Kurrent
\newline{}Versand: Stempel: »\nobreak{}\oindex{Bellinzona@\textbf{Bellinzona}|pwk}Bell{[}inzona{]}
                                       spediz. lettere, 31 VIII 2\textcolor{gray}{4}, 22\nobreak{}«.  }\pstart{}{\pb}Hrn Gustav Schwarzkopf\pend{}\pstart{}\textcolor{pink}{Wien I}\oindex{I., Innere Stadt@\textbf{I., Innere Stadt}, \emph{Verwaltungsgebiet}|pw}{}\ledrightnote{\textcolor{pink}{I., Innere Stadt}}\pend{}\pstart{}\textcolor{pink}{Tiefer Graben 17}\oindex{Wien@\textbf{Wien}!I., Innere Stadt@\textbf{I., Innere Stadt}!Tiefer Graben 17@\textbf{Tiefer Graben 17}, \emph{Wohngebäude}|pw}{}\ledrightnote{\textcolor{pink}{Tiefer Graben 17}}.\pend{}{\bigskip}
\pstart
           \noindent{}\centering{}{\pb}\textcolor{pink}{\textcolor{gray}{\textbf{Locarno}}}\oindex{Locarno@\textbf{Locarno}, \emph{Hauptstadt}|pw}{}\ledrightnote{\textcolor{pink}{Locarno}}\pend
           \vspace{1em}
\pstart
           \raggedleft{}{\pb}\textcolor{pink}{Locarno}\oindex{Locarno@\textbf{Locarno}, \emph{Hauptstadt}|pw}{}\ledrightnote{\textcolor{pink}{Locarno}}, 31. 8. 24\pend
           \vspace{0.5em}
\pstart
           Herzliche Grüße! – Nach \textcolor{pink}{Celerina}\oindex{Celerina@\textbf{Celerina}|pw}{}\ledrightnote{\textcolor{pink}{Celerina}}, über \textcolor{pink}{Luzern}\oindex{Luzern@\textbf{Luzern}|pw}{}\ledrightnote{\textcolor{pink}{Luzern}},  nach \textcolor{pink}{Lugano}\oindex{Lugano@\textbf{Lugano}, \emph{Hauptstadt}|pw}{}\ledrightnote{\textcolor{pink}{Lugano}},
                  \strikeout{von} wo allerhand Ausflüge, z. B. hieher. In
               wenigen Tagen fahr ich mit \textcolor{blue}{Lili}\pwindex{Cappellini, Lili 13.\,9.\,1909 Wien – 26.\,7.\,1928 Venedig@\textsc{Cappellini, Lili} (13.\,9.\,1909 Wien – 26.\,7.\,1928 Venedig)|pw}{}\ledrightnote{\textcolor{blue}{Lili Cappellini}} von \textcolor{pink}{Zürich}\oindex{Zürich@\textbf{Zürich}|pw}{}\ledrightnote{\textcolor{pink}{Zürich}} nach Hause\pend
           
\pstart
           {\pb}Auf baldgs Wiedersehen.{\\[\baselineskip]}Ihr{\\[\baselineskip]}\spacefill\mbox{A.}\pend
           \leftskip=0em{}\selectlanguage{ngerman}\endnumbering\briefempfaengerindex{Schwarzkopf, Gustav@\textsc{Schwarzkopf, Gustav}!zzzSchnitzler, Arthur@\emph{von Arthur Schnitzler}!1924-08-311@{31. 8. 1924}|)be}\mylabel{L04355h}
\begin{anhang}
\end{anhang}\newcommand{\dateiname}{L04355}\newcommand{\titel}{Arthur Schnitzler an Gustav Schwarzkopf, 31. 8. 1924}\newcommand{\editorInnen}{Herausgegeben von Jahnke, SelmaMüller, Martin Anton}\input{../tex-inputs/latex-korrekturansicht-abspann}
